\documentclass[UTF8,12pt]{ctexart}
\usepackage[paperwidth=240mm,paperheight=200mm,margin=14mm,top=8mm]{geometry}
\usepackage{amsmath,amssymb,mathtools}
\usepackage{xcolor}
\pagestyle{empty}
\setlength{\parindent}{0pt}
\setlength{\parskip}{0pt}
\newcommand{\qn}[1]{\textbf{\large #1}\ }
\xeCJKDeclareCharClass{CJK}{"2460 -> "24FF}
\begin{document}
\qn{2016.1}
设 $\alpha_1 = x(\cos\sqrt{x} - 1)$，$\alpha_2 = \sqrt{x}\ln(1 + \sqrt[3]{x})$，$\alpha_3 = \sqrt[3]{x + 1} - 1$. 当 $x \to 0^+$ 时，以上3个无穷小量按照从低阶到高阶的排序是（　）
\\
\noindent （A）$\alpha_1, \alpha_2, \alpha_3$.
\\
\noindent （B）$\alpha_2, \alpha_3, \alpha_1$.
\\
\noindent （C）$\alpha_2, \alpha_1, \alpha_3$.
\\
\noindent （D）$\alpha_3, \alpha_2, \alpha_1$.

\vfill
\newpage
\qn{2016.2}
已知函数 $f(x) = \begin{cases} 2(x - 1), & x < 1, \\ \ln x, & x \geqslant 1, \end{cases}$ 则 $f(x)$ 的一个原函数是（　）
\\
\noindent （A）$F(x) = \begin{cases} (x - 1)^2, & x < 1, \\ x(\ln x - 1), & x \geqslant 1. \end{cases}$
\\
\noindent （B）$F(x) = \begin{cases} (x - 1)^2, & x < 1, \\ x(\ln x + 1) - 1, & x \geqslant 1. \end{cases}$
\\
\noindent （C）$F(x) = \begin{cases} (x - 1)^2, & x < 1, \\ x(\ln x + 1) + 1, & x \geqslant 1. \end{cases}$
\\
\noindent （D）$F(x) = \begin{cases} (x - 1)^2, & x < 1, \\ x(\ln x - 1) + 1, & x \geqslant 1. \end{cases}$

\vfill
\newpage
\qn{2016.3}
反常积分 ① $\int_{-\infty}^0 \frac{1}{x^2}e^{\frac{1}{x}}\,\mathrm{d}x$，② $\int_0^{+\infty} \frac{1}{x^2}e^{\frac{1}{x}}\,\mathrm{d}x$ 的敛散性为（　）
\\
\noindent （A）①收敛，②收敛.
\\
\noindent （B）①收敛，②发散.
\\
\noindent （C）①发散，②收敛.
\\
\noindent （D）①发散，②发散.

\vfill
\newpage
\qn{2016.4}
设函数 $f(x)$ 在 $(-\infty, +\infty)$ 内连续，其导函数的图形如图所示，则（　）
\\
\noindent （A）函数 $f(x)$ 有2个极值点，曲线 $y = f(x)$ 有2个拐点.
\\
\noindent （B）函数 $f(x)$ 有2个极值点，曲线 $y = f(x)$ 有3个拐点.
\\
\noindent （C）函数 $f(x)$ 有3个极值点，曲线 $y = f(x)$ 有1个拐点.
\\
\noindent （D）函数 $f(x)$ 有3个极值点，曲线 $y = f(x)$ 有2个拐点.

\vfill
\newpage
\qn{2016.5}
设函数 $f_i(x) \ (i = 1, 2)$ 具有二阶连续导数，且 $f_i''(x_0) < 0 \ (i = 1, 2)$. 若两条曲线 $y = f_i(x) \ (i = 1, 2)$ 在点 $(x_0, y_0)$ 处具有公切线 $y = g(x)$，且在该点处曲线 $y = f_1(x)$ 的曲率大于曲线 $y = f_2(x)$ 的曲率，则在 $x_0$ 的某个邻域内，有（　）
\\
\noindent （A）$f_1(x) \leqslant f_2(x) \leqslant g(x)$.
\\
\noindent （B）$f_2(x) \leqslant f_1(x) \leqslant g(x)$.
\\
\noindent （C）$f_1(x) \leqslant g(x) \leqslant f_2(x)$.
\\
\noindent （D）$f_2(x) \leqslant g(x) \leqslant f_1(x)$.

\vfill
\newpage
\qn{2016.6}
已知函数 $f(x, y) = \frac{e^x}{x - y}$，则（　）
\\
\noindent （A）$f'_x - f'_y = 0$.
\\
\noindent （B）$f'_x + f'_y = 0$.
\\
\noindent （C）$f'_x - f'_y = f$.
\\
\noindent （D）$f'_x + f'_y = f$.

\vfill
\newpage
\qn{2016.7}
设 $\boldsymbol{A}, \boldsymbol{B}$ 是可逆矩阵，且 $\boldsymbol{A}$ 与 $\boldsymbol{B}$ 相似，则下列结论错误的是（　）
\\
\noindent （A）$\boldsymbol{A}^{\mathrm{T}}$ 与 $\boldsymbol{B}^{\mathrm{T}}$ 相似.
\\
\noindent （B）$\boldsymbol{A}^{-1}$ 与 $\boldsymbol{B}^{-1}$ 相似.
\\
\noindent （C）$\boldsymbol{A} + \boldsymbol{A}^{\mathrm{T}}$ 与 $\boldsymbol{B} + \boldsymbol{B}^{\mathrm{T}}$ 相似.
\\
\noindent （D）$\boldsymbol{A} + \boldsymbol{A}^{-1}$ 与 $\boldsymbol{B} + \boldsymbol{B}^{-1}$ 相似.

\vfill
\newpage
\qn{2016.8}
设二次型 $f(x_1, x_2, x_3) = a(x_1^2 + x_2^2 + x_3^2) + 2x_1x_2 + 2x_2x_3 + 2x_1x_3$ 的正、负惯性指数分别为1，2，则（　）
\\
\noindent （A）$a > 1$.
\\
\noindent （B）$a < -2$.
\\
\noindent （C）$-2 < a < 1$.
\\
\noindent （D）$a = 1$ 或 $a = -2$.

\vfill
\newpage
\qn{2016.9}
曲线 $y = \frac{x^3}{1 + x^2} + \arctan(1 + x^2)$ 的斜渐近线方程为 ______.

\vfill
\newpage
\qn{2016.10}
极限 $\lim_{n \to \infty} \frac{1}{n^2}\left(\sin\frac{1}{n} + 2\sin\frac{2}{n} + \cdots + n\sin\frac{n}{n}\right) =$ ______.

\vfill
\newpage
\qn{2016.11}
以 $y = x^2 - e^x$ 和 $y = x^2$ 为特解的一阶非齐次线性微分方程为 ______.

\vfill
\newpage
\qn{2016.12}
已知函数 $f(x)$ 在 $(-\infty, +\infty)$ 上连续，且 $f(x) = (x + 1)^2 + 2\int_0^x f(t)\,\mathrm{d}t$，则当 $n \geqslant 2$ 时，$f^{(n)}(0) =$ ______.

\vfill
\newpage
\qn{2016.13}
已知动点 $P$ 在曲线 $y = x^3$ 上运动，记坐标原点与点 $P$ 间的距离为 $l$. 若点 $P$ 的横坐标对时间的变化率为常数 $v_0$，则当点 $P$ 运动到点 $(1, 1)$ 时，$l$ 对时间的变化率是 ______.

\vfill
\newpage
\qn{2016.14}
设矩阵 $\boldsymbol{A} = \begin{pmatrix} a & -1 & -1 \\ -1 & a & -1 \\ -1 & -1 & a \end{pmatrix}$ 与矩阵 $\begin{pmatrix} 1 & 1 & 0 \\ 0 & -1 & 1 \\ 1 & 0 & 1 \end{pmatrix}$ 等价，则 $a =$ ______.

\vfill
\newpage
\qn{2016.15}
（本题满分10分）
求极限 $\lim_{x \to 0} (\cos 2x + 2x\sin x)^{\frac{1}{x^4}}$.

\vfill
\newpage
\qn{2016.16}
（本题满分10分）
设函数 $f(x) = \int_0^1 |t^2 - x^2|\,\mathrm{d}t \ (x > 0)$，求 $f'(x)$，并求 $f(x)$ 的最小值.

\vfill
\newpage
\qn{2016.17}
（本题满分10分）
已知函数 $z = z(x, y)$ 由方程 $(x^2 + y^2)z + \ln z + 2(x + y + 1) = 0$ 确定，求 $z = z(x, y)$ 的极值.

\vfill
\newpage
\qn{2016.18}
（本题满分10分）
设 $D$ 是由直线 $y = 1$，$y = x$，$y = -x$ 围成的有界区域，计算二重积分 $\iint_D \frac{x^2 - xy - y^2}{x^2 + y^2}\,\mathrm{d}x\mathrm{d}y$.

\vfill
\newpage
\qn{2016.19}
（本题满分10分）
已知 $y_1(x) = e^x$，$y_2(x) = u(x)e^x$ 是二阶微分方程 $(2x - 1)y'' - (2x + 1)y' + 2y = 0$ 的两个解. 若 $u(-1) = e$，$u(0) = -1$，求 $u(x)$，并写出该微分方程的通解.

\vfill
\newpage
\qn{2016.20}
（本题满分11分）
设 $D$ 是由曲线 $y = \sqrt{1 - x^2} \ (0 \leqslant x \leqslant 1)$ 与 $\begin{cases} x = \cos^3 t, \\ y = \sin^3 t, \end{cases} \ \left(0 \leqslant t \leqslant \frac{\pi}{2}\right)$ 围成的平面区域，求 $D$ 绕 $x$ 轴旋转一周所得旋转体的体积和表面积.

\vfill
\newpage
\qn{2016.21}
（本题满分11分）
已知函数 $f(x)$ 在区间 $\left[0, \frac{3\pi}{2}\right]$ 上连续，在 $\left(0, \frac{3\pi}{2}\right)$ 内是函数 $\frac{\cos x}{2x - 3\pi}$ 的一个原函数，且 $f(0) = 0$.
\\
（Ⅰ）求 $f(x)$ 在区间 $\left[0, \frac{3\pi}{2}\right]$ 上的平均值；
\\
（Ⅱ）证明 $f(x)$ 在区间 $\left(0, \frac{3\pi}{2}\right)$ 内存在唯一零点.

\vfill
\newpage
\qn{2016.22}
（本题满分11分）
设矩阵 $\boldsymbol{A} = \begin{pmatrix} 1 & 1 & 1-a \\ 1 & 0 & a \\ a+1 & 1 & a+1 \end{pmatrix}$，$\boldsymbol{\beta} = \begin{pmatrix} 0 \\ 1 \\ 2a-2 \end{pmatrix}$，且方程组 $\boldsymbol{Ax} = \boldsymbol{\beta}$ 无解.
\\
（Ⅰ）求 $a$ 的值；
\\
（Ⅱ）求方程组 $\boldsymbol{A}^{\mathrm{T}}\boldsymbol{Ax} = \boldsymbol{A}^{\mathrm{T}}\boldsymbol{\beta}$ 的通解.

\vfill
\newpage
\qn{2016.23}
（本题满分11分）
已知矩阵 $\boldsymbol{A} = \begin{pmatrix} 0 & -1 & 1 \\ 2 & -3 & 0 \\ 0 & 0 & 0 \end{pmatrix}$.
\\
（Ⅰ）求 $\boldsymbol{A}^{99}$；
\\
（Ⅱ）设3阶矩阵 $\boldsymbol{B} = (\boldsymbol{\alpha}_1, \boldsymbol{\alpha}_2, \boldsymbol{\alpha}_3)$ 满足 $\boldsymbol{B}^2 = \boldsymbol{BA}$. 记 $\boldsymbol{B}^{100} = (\boldsymbol{\beta}_1, \boldsymbol{\beta}_2, \boldsymbol{\beta}_3)$，将 $\boldsymbol{\beta}_1, \boldsymbol{\beta}_2, \boldsymbol{\beta}_3$ 分别表示为 $\boldsymbol{\alpha}_1, \boldsymbol{\alpha}_2, \boldsymbol{\alpha}_3$ 的线性组合.

\vfill
\newpage
\end{document}